%! AP Statistics — Unit 1, Lesson 5
\documentclass[10pt]{article}
\usepackage[margin=0.6in]{geometry}
\usepackage{xcolor}
\usepackage[most]{tcolorbox}
\usepackage{enumitem}
\usepackage{multicol}
\usepackage{amsmath,amssymb}
\usepackage{array}
\usepackage{tabularx}
\tcbuselibrary{breakable}

\definecolor{sky}{HTML}{EAF3FB}
\definecolor{navy}{HTML}{1F3A5F}
\definecolor{redbox}{HTML}{FCECEC}
\definecolor{greenbox}{HTML}{EDF8F0}
\definecolor{goldbox}{HTML}{FFF7E6}
\newtcolorbox{skillbox}[2][]{
    colback=#2, colframe=navy, boxrule=0.8pt, arc=2mm,
    left=2mm,right=2mm,top=1.5mm,bottom=1.5mm,
    fonttitle=\bfseries, title={#1}, halign title=center, breakable
}
\setlist[itemize]{leftmargin=1.2em,itemsep=0.35em,topsep=0.35em}
\setlist[enumerate]{leftmargin=1.4em,itemsep=0.35em,topsep=0.35em}
\newcolumntype{C}[1]{>{\centering\arraybackslash}p{#1}}
\newcolumntype{Y}{>{\centering\arraybackslash}X}

\newcommand{\CourseName}{AP Statistics}
\newcommand{\SchoolYear}{2026--2027}
\newcommand{\MeetingLength}{55 minutes}
\newcommand{\UnitNumberName}{Unit 1: Exploring One-Variable Data \quad}
\newcommand{\LessonNumberName}{Lesson 1.5: Representing a Quantitative Variable with Graphs}

\begin{document}
\begin{center}
    {\Large\bfseries \CourseName: \SchoolYear\ (\MeetingLength) \\
    {\small\bfseries \UnitNumberName \LessonNumberName}}
\end{center}

\begin{tcolorbox}[colback=sky,colframe=navy,boxrule=0.9pt,arc=2mm,
    left=2mm,right=2mm,top=1.5mm,bottom=1.5mm]
    \textbf{Primary Objective:} Students will construct and interpret dotplots, stemplots, and
    histograms for quantitative data. Students will recognize that the shape of a distribution
    (visible in these displays) is the foundation for describing and comparing data in subsequent
    lessons (Big Idea: UNC).
\end{tcolorbox}

\begin{skillbox}[Priority Ideas \& Skills]{goldbox}
\begin{minipage}[t]{0.48\linewidth}
    \textbf{AP Skill 2 --- Data Analysis}
    \begin{itemize}
        \item Choose an appropriate graph type for a given quantitative dataset.
        \item Construct a correctly labeled dotplot, stemplot, and histogram from raw data.
        \item Interpret a graph: identify approximate center, spread, and shape (formal description comes in Lesson 1.6).
    \end{itemize}
\end{minipage}
\hfill
\begin{minipage}[t]{0.48\linewidth}
    \textbf{Key Understandings}
    \begin{itemize}
        \item Histogram bars \textit{touch} because quantitative data are continuous (or ordered); categories have a gap.
        \item Dotplots preserve every value; histograms group data into intervals --- each trade-off matters.
        \item Stem-and-leaf plots show the actual data values while giving a visual shape.
    \end{itemize}
\end{minipage}
\end{skillbox}

\begin{skillbox}[{Vocabulary, Concepts, \& Theorems}]{greenbox}
\begin{minipage}[t]{0.48\linewidth}
    \begin{itemize}
        \item \textbf{Dotplot} --- each data value represented by a dot above a number line; reveals every value and clusters
        \item \textbf{Stemplot (stem-and-leaf plot)} --- stems are leading digit(s); leaves are trailing digits; preserves individual values; back-to-back for comparison
        \item \textbf{Histogram} --- data grouped into equal-width intervals (bins); bar height = frequency or relative frequency; bars touch
    \end{itemize}
\end{minipage}
\hfill
\begin{minipage}[t]{0.48\linewidth}
    \begin{itemize}
        \item \textbf{Bin (class interval)} --- the equal-width grouping used in a histogram
        \item \textbf{Shape} --- the overall pattern of the distribution (symmetric, skewed left/right, unimodal, bimodal, uniform) --- formal description in Lesson 1.6
        \item \textbf{Gap} --- a range of values with no observations; \textbf{cluster} --- a group of values that appear together
    \end{itemize}
\end{minipage}
\end{skillbox}

\begin{skillbox}[Activate Prior Knowledge \& Spiral Review]{sky}
\begin{minipage}[t]{0.48\linewidth}
    \textbf{Quick Review (3 min)}\\
    Review: what makes a variable quantitative? Why do bar chart bars NOT touch? Bridge: when data are quantitative (continuous or ordered), the bars \textit{do} touch because adjacent bins share a boundary.
\end{minipage}
\hfill
\begin{minipage}[t]{0.48\linewidth}
    \textbf{Spiral Connections}
    \begin{itemize}
        \item Lesson 1.6: use these graphs to describe shape, center, spread, and outliers (SOCS).
        \item Lesson 1.8: boxplots are another quantitative display --- compare them to histograms.
        \item Unit 4: histogram shape connects to probability distributions.
    \end{itemize}
\end{minipage}
\end{skillbox}

\begin{skillbox}[Hook \& Lesson]{sky}
\begin{minipage}[t]{0.48\linewidth}
    \textbf{Hook (4--5 min)}\\
    Collect class data: ``How many minutes did you sleep last night?'' Record on the board. Ask: ``Without calculating anything, what do you notice about this list of numbers?'' Students observe variability, a rough center, and possible extremes. Today's graphs make those observations precise.
\end{minipage}
\hfill
\begin{minipage}[t]{0.48\linewidth}
    \textbf{Lesson Flow (15 min)}
    \begin{enumerate}
        \item Build a dotplot of the sleep data together.
        \item Rearrange the same data into a stemplot; compare what's visible.
        \item Group the data into bins and construct a histogram; choose bin width together.
        \item Discuss trade-offs: dotplot (all values, small $n$), stemplot (sorted values), histogram (large $n$, grouped).
    \end{enumerate}
\end{minipage}
\end{skillbox}

\begin{skillbox}[Explicit Instruction \& Active Monitoring]{sky}
\begin{minipage}[t]{0.48\linewidth}
    \textbf{Direct Instruction (10 min)}
    \begin{itemize}
        \item Histogram construction rules: equal bin widths, bins start where the last ends, label axes with variable and scale, add a title.
        \item Demonstrate how changing bin width changes the histogram's appearance --- discuss the judgment involved.
        \item Stemplot rule: split stems if needed; back-to-back stemplot preview for Lesson 1.9.
    \end{itemize}
\end{minipage}
\hfill
\begin{minipage}[t]{0.48\linewidth}
    \textbf{Active Monitoring}
    \begin{itemize}
        \item Circulate as students construct graphs; watch for: bars with gaps in histograms, unlabeled axes, unequal bin widths.
        \item Ask students: ``If you doubled the bin width, what would happen to the shape of your histogram?''
        \item Probe: ``Which graph would you use if you only had 10 data values? 200?''
    \end{itemize}
\end{minipage}
\end{skillbox}

\begin{skillbox}[Group Work \& Differentiation]{redbox}
\begin{minipage}[t]{0.48\linewidth}
    \textbf{Group Activity (8--10 min)}\\
    Groups receive a dataset of 25 test scores (0--100). Each group constructs:
    \begin{enumerate}
        \item A dotplot
        \item A stemplot (stems: tens digit; leaves: units digit)
        \item A histogram with bins of width 10
    \end{enumerate}
    Then groups compare: ``Which graph is most useful for a teacher? Justify.''
\end{minipage}
\hfill
\begin{minipage}[t]{0.48\linewidth}
    \textbf{Differentiation}
    \begin{itemize}
        \item \textit{Support}: Provide a dotplot template with a pre-drawn number line; stemplot with stems pre-filled.
        \item \textit{Extension}: Students experiment with bin widths of 5, 10, and 20 for the histogram. Write: how does bin width affect the story told by the graph?
        \item \textit{Technology}: Use a graphing calculator or Desmos to verify their hand-drawn histogram.
    \end{itemize}
\end{minipage}
\end{skillbox}

\begin{skillbox}[Individual Work \& Assessment]{redbox}
\begin{minipage}[t]{0.48\linewidth}
    \textbf{Exit Ticket (5 min)}\\
    Given: 12 data values (e.g., number of books read last year: 3, 7, 1, 5, 2, 8, 4, 6, 3, 10, 2, 5). Students sketch a dotplot and a histogram (bins: 0--3, 4--7, 8--11). They label both graphs completely.
\end{minipage}
\hfill
\begin{minipage}[t]{0.48\linewidth}
    \textbf{AP-Style Check}\\
    \textit{``A histogram of daily high temperatures is shown. Describe what the histogram tells you about the distribution of temperatures.''} \\[4pt]
    Students should mention: approximate center, spread (range of values visible), and overall shape --- this primes Lesson 1.6 SOCS framework.
\end{minipage}
\end{skillbox}

\begin{skillbox}[Reinforcement \& Extension]{goldbox}
\begin{minipage}[t]{0.48\linewidth}
    \textbf{Homework / Practice}
    \begin{itemize}
        \item Construct a dotplot, stemplot, and histogram from a provided dataset of 20 values.
        \item Answer: ``Which display would you choose to present to an audience unfamiliar with statistics? Why?''
    \end{itemize}
\end{minipage}
\hfill
\begin{minipage}[t]{0.48\linewidth}
    \textbf{Extension / Preview}
    \begin{itemize}
        \item \textit{Extension}: Collect real data (e.g., prices of 20 items at a grocery store) and create all three displays.
        \item \textit{Preview}: Lesson 1.6 teaches how to \textit{describe} a distribution using SOCS --- Shape, Outliers, Center, Spread. Review the graphs you made today and try to describe each one in words before next class.
    \end{itemize}
\end{minipage}
\end{skillbox}

\end{document}
